\hypersetup{pdftitle={Feature Showcase}}
\SetFooter{SAMPLE NOTES}{\begin{tabular}[t]{@{}r@{}}Feature Showcase\\Sample Institute of Technology\end{tabular}}

\begin{document}

% ------------------------------------------------------------ cover (opt-in extra)
\thispagestyle{empty}
\vspace*{\fill}
\begin{center}
  {\fontsize{34}{42}\selectfont\bfseries\clTitleColor FEATURE\\[6pt] SHOWCASE\par}
  \vspace{14pt}
  {\Large Limits, Programs and Diagrams\par}
\end{center}
\vspace*{\fill}
\clearpage
\setcounter{page}{1}

% ------------------------------------------------------------ title block
\TitleBlock{%
  {\Large\bfseries Sample Institute of Technology}\\
  {\large (An Example Institution)}\\
  {\large AY: 2026--27 (Odd Semester)}\\
  {\large\bfseries 101BS: Mathematics - I}\\[0.6em]
  {\large\bfseries Tutorial -- 1}%
}

\section*{Indeterminate Forms}

\begin{defbox}
\definition{If $\lim\limits_{x\to a} f(x)=0$ and $\lim\limits_{x\to a} g(x)=0$, the limit of
$\frac{f(x)}{g(x)}$ is called an \keyterm{indeterminate form} of the type $\frac{0}{0}$.}
\end{defbox}

Other indeterminate forms are
\[
  \frac{\infty}{\infty},\qquad 0\times\infty,\qquad \infty-\infty,\qquad 1^{\infty},\qquad
  0^{0},\qquad \infty^{0}.
\]

\subsection*{L'Hospital's Rule}

\begin{formulabox}
\begin{equation}
  \lim_{x\to a}\frac{f(x)}{g(x)}=\lim_{x\to a}\frac{f'(x)}{g'(x)} \tag{1}
\end{equation}
\end{formulabox}
\important{The rule applies only when the latter limit exists.}

\begin{keybox}
\important{Differentiate the numerator and the denominator separately,} not the quotient.
\end{keybox}

\subsection*{Worked example}

Evaluate $\displaystyle\lim_{x\to 0}\frac{3x^{2}+6\cos x-6}{x^{4}}$
\begin{align*}
  \lim_{x\to 0}\frac{3x^{2}+6\cos x-6}{x^{4}}
    &=\lim_{x\to 0}\frac{6x-6\sin x}{4x^{3}} &&\text{(0/0 form)}\\
    &=\lim_{x\to 0}\frac{6-6\cos x}{12x^{2}} &&\text{(0/0 form)}\\
    &=\lim_{x\to 0}\frac{6\sin x}{24x}
\end{align*}
\begin{answerbox}
\[ \lim_{x\to 0}\frac{3x^{2}+6\cos x-6}{x^{4}}=\frac{1}{4} \]
\end{answerbox}

\subsection*{Standard expansions}
\[
  \text{(I)}\quad e^{x}=1+x+\frac{x^{2}}{2!}+\frac{x^{3}}{3!}+\cdots
  =\sum_{n=0}^{\infty}\frac{x^{n}}{n!},\qquad(-\infty,\infty)
\]
\[
  \text{(II)}\quad \sin x=x-\frac{x^{3}}{3!}+\frac{x^{5}}{5!}-\cdots
  =\sum_{n=0}^{\infty}(-1)^{n}\frac{x^{2n+1}}{(2n+1)!},\qquad(-\infty,\infty)
\]

% ------------------------------------------------------------ table + quick reference
\section*{Greek Alphabet}
\begin{center}
\begin{tabular}{|c|c|c|c|}
  \hline
  \TableHead Symbol & Name & Symbol & Name \\ \hline
  $\alpha$ & Alpha & $\nu$ & Nu \\
  $\beta$ & Beta & $\xi$ & Xi \\
  $\gamma,\Gamma$ & Gamma & $\pi$ & Pi \\
  $\delta,\Delta$ & Delta & $\rho$ & Rho \\
  $\theta,\Theta$ & Theta & $\sigma,\Sigma$ & Sigma \\
  \hline
\end{tabular}
\end{center}

\section*{Formulas}
\begin{multicols}{2}
\subsection*{Trigonometry}
\begin{gather*}
  \sin^{2}x+\cos^{2}x=1\\
  \tan x=\frac{\sin x}{\cos x}
\end{gather*}
\subsection*{Hyperbolic}
\begin{gather*}
  \cosh^{2}x-\sinh^{2}x=1\\
  \cosh x+\sinh x=e^{x}
\end{gather*}
\end{multicols}

% ------------------------------------------------------------ program
\clearpage
\section*{Programming in C}
\noindent\clLabel{4.6} Write a C program to calculate the sum of digits of a number.

\noindent\clLabel{Solution:}
\begin{lstlisting}[language=C]
#include <stdio.h>
#include <conio.h>

void main()
{
    int n, r, sum = 0;

    clrscr();
    printf("Enter a number: ");
    scanf("%d", &n);

    while (n > 0)
    {
        r = n % 10;
        sum = sum + r;
        n = n / 10;
    }

    printf("Sum of digits = %d", sum);
    getch();
}
\end{lstlisting}
\noindent\clLabel{Output:}
\begin{lstlisting}[style=srcoutput]
Enter a number: 12345
Sum of digits = 15
\end{lstlisting}

\subsection*{Pseudocode}
\begin{lstlisting}[style=pseudocode]
PROCEDURE SumOfDigits(n)
    sum = 0
    WHILE n > 0 DO
        sum = sum + n MOD 10
        n = n DIV 10
    ENDWHILE
    RETURN sum
END
\end{lstlisting}

% ------------------------------------------------------------ flowchart
\clearpage
\noindent\clLabel{4.6 Flowchart}
\par\vspace{12pt}
\begin{center}
\begin{adjustbox}{max width=\textwidth,max totalheight=0.8\textheight}
\input{fc_showcase.tex}
\end{adjustbox}
\end{center}

% ------------------------------------------------------------ circuit, graph, chemistry
\clearpage
\section*{RC Circuit}
\begin{center}
\begin{circuitikz}[american]
  \draw (0,0) to[V, v=$V_s$] (0,3)
        to[R, l={$R_1 = 10\,\mathrm{k}\Omega$}] (4,3)
        to[C, l={$C = 100\,\mu\mathrm{F}$}] (4,0) -- (0,0);
  \draw (0,0) node[ground]{};
\end{circuitikz}
\end{center}
\important{The time constant of the circuit is} \[ \tau = RC \]

\section*{Graph of $y=\sin x$}
\begin{center}
\begin{tikzpicture}
  \begin{axis}[axis lines=middle, xlabel={$x$}, ylabel={$y$},
      xmin=-3.6, xmax=3.6, ymin=-1.4, ymax=1.4,
      xtick={-3.1416,0,3.1416}, xticklabels={$-\pi$,$0$,$\pi$},
      width=0.75\textwidth, height=5cm]
    \addplot[domain=-3.1416:3.1416, samples=200, thick, color=clSec] {sin(deg(x))};
  \end{axis}
\end{tikzpicture}
\end{center}

\section*{Chemistry}
\[ \ce{2H2 + O2 -> 2H2O} \qquad\qquad \ce{CaCO3 ->[$\Delta$] CaO + CO2 ^} \]

% ------------------------------------------------------------ questions
\section*{Q-1 Do as directed: [A]}
\begin{enumerate}[label=\arabic*.]
  \item Evaluate $\displaystyle\lim_{x\to 0}\frac{e^{x}-1}{x}$. \hfill (Winter 2012) [3]
  \item Expand $\sin\left(\dfrac{\pi}{4}+x\right)$ in powers of $x$. \hfill [4]
  \item Find the approximate value of $\sqrt[3]{9.12}$. \hfill (Summer 2019) [3]
\end{enumerate}

\end{document}
